\begin{ledgroupsized}[r]{120mm}
\footnotesize
\pstart
\noindent\textbf{Überlieferung:}
\pend
\end{ledgroupsized}
\begin{ledgroupsized}[r]{114mm}
\footnotesize
\pstart \parindent -6mm
\makebox[6mm][l]{\textit{L}}%
Aufzeichnung:
LH~IV~1, 4d, 2~Bl.~4. 
Ein Zettel (ca.~10~x~7~cm.);
unterer Rand nach Abriss schief, die übrigen Ränder beschnitten.
N.\;\ref{39763} nimmt ein Drittel der Textfläche von 
Bl.~\href{\Scan{LH_04_01_04d_02_004v}}%
{4~v\textsuperscript{o}}
ein.
Außerdem überliefert der Träger zwei weitere Leibniztexte.
Auf Bl.~\href{\Scan{LH_04_01_04d_02_004v}}%
{4~v\textsuperscript{o}},
senkrecht zur Schreibrichtung, ist folgender Briefentwurf %
fragmentarisch erhalten (Konzept zu einem nicht erhaltenen Schreiben an
\indexN{\textls{Brosseau}, Christophe 1630\textendash1717}Christophe Brosseau, siehe \cite{02148}\textit{LSB} I, 9 N.\;276):
\textlangle\textendash\textrangle\textit{avois prié Monsi}\textlangle\textit{eur}\textrangle\ %
\lbrack/\rbrack\
\textlangle \textit{ter}\textrangle\textit{mes exprés de me}\textlangle\textendash\textrangle\ %
\lbrack/\rbrack\
\textlangle\textit{ga}\textrangle\textit{zettes sur mon comp}\textlangle\textit{te}\textrangle\
\lbrack/\rbrack\
\textlangle\textendash\textrangle\textit{s conviendra pas.}
\lbrack/\rbrack\
\textit{payer icy, à fin}
\lbrack/\rbrack\
\textit{comptes de S.\,A.\,E.}
\lbrack/\rbrack\
\textlangle\textendash\textrangle\textit{vint chargés}
\bracks{\textit{Text bricht ab.}} %
\textit{S.\,A.\,E.} ist eine Verbesserung von \textit{S.\,A.\,S.}
Bl.~\href{\Scan{LH_04_01_04d_02_004r}}%
{4~r\textsuperscript{o}}
überliefert eine Aufzeichnung über Descartes %
(\textit{Mons.~des Cartes a fait comme les charlatans}...), die in 
\cite{02150}\textsc{Foucher de Careil}, \textit{Nouvelles Lettres et Opuscules Inédits de Leibniz}, Paris 1857 
gedruckt vorliegt (S.\;12) und vsl. in einem Band der Sechsten Reihe erscheinen wird.
Weitere Textträger, die ursprünglich mit dem Zettel zusammenhingen, konnten noch nicht ermittelt werden. %
\pend
\end{ledgroupsized}
\selectlanguage{latin}
\vspace{8mm}
\pstart%
\normalsize%
\noindent
\edlabel{04_01_04d_02_004v_2a}%
\edtext{\href{\Scan{LH_04_01_04d_02_004v}}%
{\lbrack4~v\textsuperscript{o}\rbrack}\
$a=2b$.}{%
\lemma{\href{\Scan{LH_04_01_04d_02_004v}}%
{\lbrack4~v\textsuperscript{o}\rbrack}}%
\Bfootnote{%
\textit{(1)}~$\ind{3}A\ind{3}D^2+r$ %
\textit{(2)}~$a=2b$.~\textit{L}}}
$r=1$.
\pend
\pstart
$+ 2\sqrt{4+4}=4\sqrt{2}$
\pend
\pstart
$-ra^2b = -4$
\pend
\pstart
$4+4+2 = 8$
\pend
\pstart
\rule[-2.4ex]{0pt}{6.2ex}%
fit $\dfrac{\surd2 - 1}{2}$.%
\edlabel{04_01_04d_02_004v_2b}%
\pend
\pstart
\edlabel{04_01_04d_02_004v_1a}%
\edtext{}{%
\xxref%
{04_01_04d_02_004v_1a}{04_01_04d_02_004v_1b}%
\lemma{4.1}%
\Bfootnote{%
\textit{(1)}~2.1 %
\textit{(2)}~4.1 %
\textit{(3)}~$2a$ %
\textit{(4)}~2.2 %
\textit{(5)}~2.4 + 2
\textit{(a)}~\textit{a}
\textit{(b)}~\textit{b}
\textit{(6)}~$1.2+1.2$%
~\textit{L}}}%
$4.1$
\pend
\begin{edfigure}[0.25em]%
\includegraphics[scale=1]{%
gesamttex/edit_VIII-4/images/LH_04_01_04d_02_004v_d.pdf}
\edcaption[0em]{\lbrack\textit{Fig.~1}\rbrack}
\end{edfigure}
\pstart
$1.2+1.2$%
\edlabel{04_01_04d_02_004v_1b}%
\pend
\pstart
\rule[-2.4ex]{0pt}{6.0ex}%
$4.4b = a$.\enskip
$b=\dfrac{1}{4}$.%
\edlabel{04_01_04d_02_004v_3}%
\pend
\pstart
$\surd2 - 1$.\enskip
$16+8$.
\pend
